\documentclass[a4paper]{article}     
\usepackage{amsfonts}
\usepackage{amsmath}
\usepackage{amssymb}
\usepackage{graphicx}
\usepackage{color}

\newcommand{\Q}{\mathbb{Q}}
\newcommand{\R}{\mathbb{R}}
\newcommand{\llim}{\lim\limits}
\newcommand{\dint}{\displaystyle\int}

\begin{document}

\noindent \large Auteur: Ilyas Sefiani \\
\noindent \large Studierichting: Informatica\\
\noindent \large Oefeningenreeks 6B nr. 17.4\\

\medskip

\normalsize

$$1 \cdot 3 \cdot 5 \cdot + 2 \cdot 4 \cdot 6 + 3 \cdot 5 \cdot 7 + ... + n(n+2)(n+4)$$\\ 

$s_{n}=\displaystyle\sum\limits_{i=1}^{n} i(i+2)(i+4)= \displaystyle\sum\limits_{i=1}^{n} (i^2+2i)(i+4)$\\

$=\displaystyle\sum\limits_{i=1}^{n} (i^3 +6i^2 +8i) = \displaystyle\sum\limits_{i=1}^{n} i^3 + 6\displaystyle\sum\limits_{i=1}^{n}i^2 +8\displaystyle\sum\limits_{i=1}^{n}i$\\

Omzetten:\\

$= \dfrac{n^2(n+1)^2}{4} + \dfrac{6n(n+1)(2n+1)}{6} + \dfrac{8n(n+1)}{2}$\\

$= \dfrac{n^2(n+1)^2 + 4n(n+1)(2n+1) + 16n(n+1)}{4}$\\

$= n(n+1)\left(\dfrac{n(n+1) + 4(2n+1) + 16}{4}\right)$\\

$= n(n+1)\left(\dfrac{n^2+n + 8n+4 + 16}{4}\right)$\\

$= n(n+1)\left(\dfrac{n^2 + 9n +20}{4}\right)$\\

Teller ontbinden in factoren:\\

$= \dfrac{n(n+1)(n+4)(n+5)}{4}$\\

Bewijzen dat de algemene term van $s_{n}$ altijd een element van $\mathbb{N}$ is:\\

Zowel $n(n+1)$ als $(n+4)(n+5)$ zijn twee opeenvolgende cijfers:\\

$\Rightarrow$ deelbaar door 2 $\Rightarrow$ allemaal samen deelbaar door 4!\\

\begin{thebibliography}{999}
\end{thebibliography}
\end{document} 